%%%%%%%%%%%%%%%%%%%%%%%%%%%%%%%%%%%%%%%%%%%%%%%%%%%%%%%%%%%%%%%%%%%%%%%%%%% 
% 
% Generic template for TFC/TFM/TFG/Tesis
% 
% By:
%  + Javier Macías-Guarasa.
%    Departamento de Electrónica
%    Universidad de Alcalá
%  + Roberto Barra-Chicote.
%    Departamento de Ingeniería Electrónica
%    Universidad Politécnica de Madrid
% 
% By: + Javier Macías-Guarasa. Departamento de Electrónica Universidad de Alcalá + Roberto Barra-Chicote. Departamento de Ingeniería Electrónica Universidad Politécnica de Madrid
% 
% Based on original sources by Roberto Barra, Manuel Ocaña, Jesús Nuevo, Pedro Revenga, Fernando Herránz and Noelia Hernández. Thanks a lot to all of them, and to the many anonymous contributors found (thanks to google) that provided help in setting all this up.
% 
% See also the additionalContributors.txt file to check the name of additional contributors to this work.
% 
% If you think you can add pieces of relevant/useful examples, improvements, please contact us at (macias@depeca.uah.es)
% 
% You can freely use this template and please contribute with comments or suggestions!!!
% 
%%%%%%%%%%%%%%%%%%%%%%%%%%%%%%%%%%%%%%%%%%%%%%%%%%%%%%%%%%%%%%%%%%%%%%%%%%% 

\chapter{Introducción}
\label{cha:introduccion}

La salud mental se ha convertido en uno de los principales desafíos del siglo XXI, con más de mil millones de personas afectadas por algún trastorno mental según estimaciones recientes de la Organización Mundial de la Salud \cite{noauthor_over_nodate}. Entre los trastornos más prevalentes se encuentran la ansiedad, la depresión y los \ac{TCA}, los cuales generan un profundo impacto personal, social y económico. La identificación de estos trastornos, especialmente en etapas tempranas, continúa siendo un desafío, en el que los recientes avances tecnológicos y computacionales pueden desempeñar un papel relevante.

En paralelo, la inteligencia artificial ha experimentado un crecimiento notable en los últimos años, impulsando avances significativos en múltiples ámbitos. En la actualidad, estas tecnologías forman parte del uso cotidiano de millones de personas, siendo especialmente destacables los sistemas conversacionales (\textit{chatbots}) como \textit{ChatGPT}, que han popularizado el acceso a modelos avanzados de lenguaje. Este contexto ha dado lugar a nuevos usos no previstos inicialmente, observándose incluso casos en los que estos sistemas son utilizados como sustitutos informales de apoyo psicológico.

Esta tendencia plantea diversas cuestiones sobre el papel que pueden desempeñar estos modelos en dicho ámbito. Por un lado, no están diseñados ni validados para sustituir la evaluación y el tratamiento por parte de profesionales cualificados, lo que limita su aplicación en contextos clínicos. Por otro lado, su capacidad para analizar lenguaje y detectar patrones abre la posibilidad de que puedan resultar útiles como herramientas de apoyo en tareas específicas, como la identificación de indicadores asociados a distintos trastornos de salud mental.

Dentro de este contexto, el \ac{PLN} se ha consolidado como una de las áreas más relevantes, permitiendo el desarrollo de modelos capaces de comprender y analizar texto de forma cada vez más precisa. Modelos basados en arquitecturas \textit{transformer}, como \ac{BERT}, así como los recientes \acp{LLM}, han demostrado un alto rendimiento en tareas como la clasificación de texto, el análisis de sentimientos o la detección de emociones. Estas capacidades han favorecido su aplicación en distintos dominios, incluido el ámbito de la salud mental, donde el lenguaje puede actuar como un indicador del estado psicológico de una persona.

Por otro lado, el uso generalizado de plataformas digitales, redes sociales y servicios de comunicación ha dado lugar a un crecimiento masivo de datos textuales generados por los usuarios. Esta disponibilidad de información supone un recurso de gran valor para los sistemas de inteligencia artificial, ya que permite entrenar y evaluar modelos en contextos reales. En el ámbito de la salud mental, estos datos pueden contener señales lingüísticas asociadas a estados emocionales o patrones de comportamiento, lo que abre la posibilidad de desarrollar herramientas automáticas capaces de identificar indicios tempranos de determinados trastornos a partir del análisis del lenguaje.

En este contexto, el presente trabajo tiene como objetivo evaluar y comparar la capacidad de distintos modelos de procesamiento del lenguaje natural para la detección automática de indicadores de salud mental en textos escritos en español, centrándose en ansiedad, depresión y \ac{TCA}. Para ello, se realiza una comparación entre modelos encoder tradicionales basados en \ac{BERT}, \acp{LLM} open-source y \acp{LLM} propietarios de última generación.

Con el fin de alcanzar este objetivo general, se plantean los siguientes objetivos específicos:

\begin{itemize}
	\item Analizar y preprocesar un conjunto de datos en español compuesto por textos extraídos de redes sociales, que contenga anotaciones sobre la presencia o ausencia de ansiedad, depresión y trastornos de la conducta alimentaria.
	
	\item Plantear tanto tareas de clasificación binaria independientes para ansiedad, depresión y \ac{TCA}, como una tarea de clasificación multiclase. Las tareas binarias permiten evaluar la capacidad de los modelos para detectar la presencia o ausencia de cada trastorno de forma individual, mientras que la tarea multiclase permite analizar su capacidad para distinguir simultáneamente las clases control, ansiedad, depresión y \ac{TCA}, así como identificar cuales de ellas tienden a confundirse entre sí.
	
	\item Entrenar y evaluar modelos encoder basados en \ac{BERT} mediante técnicas de \textit{fine-tuning}, incluyendo arquitecturas preentrenadas tanto específicamente para español como modelos multilingües y de contexto largo.
	
	\item Aplicar \acp{LLM}, tanto open-source como propietarios, mediante estrategias de \textit{zero-shot prompting} a través de \acp{API} de distintos proveedores, con el fin de realizar predicciones sin necesidad de reentrenamiento supervisado.
	
	\item Comparar de forma rigurosa ambos enfoques mediante métricas de rendimiento y estabilidad, análisis de errores, técnicas de comparación estadística entre modelos y evaluación de la relación entre rendimiento y coste.
\end{itemize}

Finalmente, cabe destacar que este trabajo se enmarca en el ámbito del cribado automático y el apoyo al análisis preliminar. En ningún caso pretende sustituir el diagnóstico clínico ni la evaluación realizada por profesionales de la salud mental, sino contribuir al desarrollo de herramientas que puedan facilitar la identificación temprana de posibles señales de riesgo.

%%% Local Variables:
%%% TeX-master: "../book"
%%% End:


